\section{4.9.Integración y
distribución}\label{integraciuxf3n-y-distribuciuxf3n}

\subsection{4.9. Distribución e integración de la
solución}\label{distribuciuxf3n-e-integraciuxf3n-de-la-soluciuxf3n}

Uno de los objetivos perseguidos durante el desarrollo de EVMAudit fue
diseñar una solución desacoplada de cualquier interfaz concreta y
suficientemente flexible para poder ser reutilizada en distintos
contextos.

Por este motivo, la lógica principal de análisis se implementó como una
librería independiente, separando completamente las funcionalidades
relacionadas con la detección y procesamiento de vulnerabilidades de los
mecanismos de interacción con el usuario.

Esta decisión permitió demostrar la reutilización de la solución
desarrollada mediante su integración en otros componentes software sin
necesidad de modificar el núcleo de la herramienta.

\subsubsection{4.9.1. Distribución de la
librería}\label{distribuciuxf3n-de-la-libreruxeda}

La implementación de EVMAudit como una librería independiente permite
desacoplar el motor de análisis de cualquier entorno concreto de
ejecución.

Este enfoque proporciona diversas ventajas:

\begin{itemize}
\tightlist
\item
  reutilización de la herramienta desde otros proyectos
\item
  separación entre lógica de negocio e interfaces de usuario
\item
  facilidad de mantenimiento
\item
  posibilidad de incorporar nuevas interfaces en el futuro.
\end{itemize}

La distribución de la librería permite que el proceso de análisis pueda
ejecutarse desde aplicaciones externas sin necesidad de duplicar
funcionalidades.

!TODO: referencia cruzada al Anexo A (publicación en PyPI).

\subsubsection{4.9.2. Aplicación web como caso de
uso}\label{aplicaciuxf3n-web-como-caso-de-uso}

Con el objetivo de validar la reutilización de la arquitectura
desarrollada, se implementó una aplicación web que utiliza EVMAudit como
motor de análisis.

La aplicación no implementa mecanismos propios de detección de
vulnerabilidades, sino que delega completamente las tareas de análisis
en la librería desarrollada.

Esta aproximación permite mantener una única implementación del proceso
de análisis y garantiza la consistencia de los resultados obtenidos.

La aplicación actúa, por tanto, como una capa de presentación encargada
de:

\begin{itemize}
\tightlist
\item
  recibir contratos inteligentes proporcionados por el usuario
\item
  invocar las funciones de EVMAudit
\item
  mostrar el progreso del análisis
\item
  presentar los resultados obtenidos.
\end{itemize}

De esta forma, la aplicación web constituye una demostración práctica de
la reutilización de la solución desarrollada.

\subsubsection{4.9.3. Flujo de funcionamiento de la
aplicación}\label{flujo-de-funcionamiento-de-la-aplicaciuxf3n}

La aplicación desarrollada permite dos mecanismos de entrada:

\begin{itemize}
\tightlist
\item
  carga de contratos Solidity mediante fichero
\item
  introducción directa del código fuente a través de la interfaz web.
\end{itemize}

Una vez recibido el contrato, la aplicación inicia el proceso de
análisis y ejecuta internamente el pipeline implementado por EVMAudit.

Las distintas fases del proceso se ejecutan de forma secuencial:

\begin{enumerate}
\def\labelenumi{\arabic{enumi}.}
\tightlist
\item
  ejecución de Slither
\item
  ejecución de Mythril
\item
  normalización de resultados
\item
  correlación de vulnerabilidades
\item
  generación de propiedades
\item
  ejecución de Echidna
\item
  generación del informe final.
\end{enumerate}

Durante la ejecución, la aplicación informa al usuario del progreso
alcanzado y permite recuperar el informe una vez finalizado el análisis.

Este comportamiento demuestra que la arquitectura modular adoptada
permite integrar EVMAudit en aplicaciones externas sin modificar la
lógica interna del sistema.

!TODO: insertar diagrama simplificado de la aplicación web.

!TODO: insertar captura de pantalla de la interfaz.

!TODO: referencia cruzada al Anexo B (aplicación web).

\subsubsection{4.9.4. Separación entre presentación y lógica de
análisis}\label{separaciuxf3n-entre-presentaciuxf3n-y-luxf3gica-de-anuxe1lisis}

La decisión de implementar EVMAudit como una librería independiente
permite mantener separadas las responsabilidades del sistema.

Mientras que la librería se encarga de:

\begin{itemize}
\tightlist
\item
  ejecutar herramientas externas
\item
  procesar resultados
\item
  correlacionar vulnerabilidades
\item
  generar informes
\end{itemize}

la aplicación web únicamente proporciona los mecanismos de interacción
con el usuario.

Esta separación favorece:

\begin{itemize}
\tightlist
\item
  la mantenibilidad del sistema
\item
  la reutilización de la solución
\item
  la incorporación de nuevas interfaces
\item
  la evolución independiente de cada componente.
\end{itemize}

En consecuencia, la aplicación web debe entenderse como una demostración
de la capacidad de integración de EVMAudit y no como la contribución
principal del trabajo.

\subsubsection{4.9.5. Consideraciones sobre
infraestructura}\label{consideraciones-sobre-infraestructura}

Los aspectos relacionados con la publicación de la librería, el
despliegue de la aplicación y la infraestructura utilizada no
constituyen el núcleo de la contribución desarrollada en este Trabajo
Fin de Máster.

Por este motivo, dichos elementos se incluyen en los anexos del
documento.

!TODO: referencia cruzada al Anexo A (PyPI).

!TODO: referencia cruzada al Anexo B (Aplicación web).

!TODO: referencia cruzada al Anexo C (Docker).

!TODO: referencia cruzada al Anexo D (Despliegue).
